% ============================================================
%  MODULO 19 — Podcast: apresentacao consolidada em portugues
%  SPEC R683/R684 — faixa principal PT garantida + episodio NotebookLM
% ============================================================
\chapter{Podcast: O Livro Inteiro em Uma Explica\c{c}\~ao \'{U}nica}

\roteirodidatico
{Este capitulo e o livro falado: o que e o Core, o consolidado de todas as partes, o especial, o limite, a ajuda ao pesquisador e o funcionamento com minucia e justificativa, num equilibrio unico entre leigo e dev.}
{\emph{Apresenta\c{c}\~ao} diz o que e; \emph{consolidado} resume o todo; \emph{especial} mostra a vantagem; \emph{limite} mostra a fronteira; \emph{minucia} abre o mecanismo com conta e motivo.}
{Ou\c{c}a a faixa principal com o roteiro na m\~ao, pause em cada bloco e confira o termo no dicionario nivel 0.}
{Audio resume; prova esta no livro e nas fontes com DOI. Voz gerada pode errar enfase.}
{Que bloco voce reconta sozinho: a casa, os sete passos, o especial ou o limite?}

\casoguiado{Ana e Bruno ouvem juntos}
{Ana ouve a casa e os sete passos e anota: sem fonte, devolve. Bruno ouve Jaccard e PPV e anota as contas para refazer. Os dois terminam sabendo pedir e conferir, cada um no seu nivel.}

\section{As duas faixas e os arquivos}

\elemento{Podcast consolidado em portugues}{\metainfo{ID}{POD-02} \hfill \metainfo{Idioma}{pt-BR} \hfill \metainfo{Pasta}{output/podcast/}}

\begin{descricao}
Faixa principal em portugues, \pth{livro-core-apresentacao-ptbr.mp3}, cerca de quatro minutos, lida do roteiro consolidado \pth{roteiro-consolidado-ptbr.md}: apresenta\c{c}\~ao, consolidado em seis frases, especial, limites, ajuda ao pesquisador e minucia com Jaccard e PPV. Episodio longo do NotebookLM em pt-BR, \pth{livro-core-notebooklm-ptbr.m4a}, cerca de vinte minutos, gerado no caderno em portugues com idioma forcado a partir das quatro fontes em portugues. Apoio: sintese do caderno, nove cartoes e feed com tamanho real. Para refazer a faixa PT, rode o gerador de voz sobre o roteiro; para o episodio, crie o audio no caderno com idioma pt-BR e foco consolidado, depois baixe o audio.
\end{descricao}

\begin{funcao}
Uma explica\c{c}\~ao unica com dois ritmos: quatro minutos para situar-se, vinte para aprofundar, ambos terminando no mesmo dicionario e na mesma planta de doze passos.
\end{funcao}

\section{O consolidado em seis blocos para recontar}

\begin{descricao}
Um, o que e: casa com porteiro unico, 378 elementos e regra dura de combinado e teste. Dois, o livro em seis frases: prologo, quatro esta\c{c}\~oes, oito camadas, sete passos, memoria que aprende, qualidade e reprodu\c{c}\~ao. Tres, o especial: balcao que coordena, lousa voluntaria, caderno que vira erro em criterio, fiscal e proveni\^encia. Quatro, o limite: nota e parecido, teste e criterio, sem selo sem externa. Cinco, o pesquisador: fontes, resumo sem inventar, metodo e auditoria ate o DOI. Seis, a minucia: E1 a E7 com Jaccard e aten\c{c}\~ao, risco de 89 por cento em 216 saidas e PPV de 62 por cento em campo dificil, conferidos com doctor e recompila\c{c}\~ao.
\end{descricao}

\begin{aviso}
Limite do audio. Em divergencia entre ouvido e impresso, vale o impresso com a fonte citada. Nenhuma faixa promete sotaque, trilha ou dura\c{c}\~ao exata em outro gerador.
\end{aviso}
